\section{Introduction}

County-level outage records reveal the fraction of customers without service, not the underlying counts of newly interrupted and newly restored customers. Utility records can retrospectively decompose resilience curves into overlapping outage and restoration processes \citep{carrington2021}; in contrast, a prominent statistical outage model treats the simultaneous movements through a single net process because separate modeling is considered impractical \citep{zhu2021}. These positions pose a statistical question that is more precise than whether two processes can be named: when does their separation improve prediction from aggregate state observations?

We study the bounded-state transition
\begin{equation}
Y_{t+1}=Y_t+U_0(X_t)(1-Y_t)-R_0(X_t)Y_t+\varepsilon_t,
\label{eq:intro_transition}
\end{equation}
where $U_0$ is an interruption component, $R_0$ is a restoration component, and $X_t$ denotes observed drivers. The competing one-flow representation uses one signed function: positive values act on the served pool $1-Y_t$, negative values act on the interrupted pool $Y_t$. It can switch direction across driver values, but it cannot keep both directions positive at the same driver value. The two-flow representation is therefore more expressive, but it also estimates an additional function. The relevant question is not whether two flows are possible; it is whether their representation gain repays their estimation cost.

The analysis has one organizing quantity. In a fixed-design Gaussian benchmark, let $G_n$ denote the approximation loss of the best one-flow span and let $\sigma_\varepsilon^2$ be the one-step noise variance. The ratio $\Gamma_n=nG_n/\sigma_\varepsilon^2$ factors into information for total turnover and the cost of suppressing one active direction. This factorization unifies three issues that are often discussed separately: whether two flows are needed, whether they can be identified, and whether the additional component is worth estimating.

Our contributions follow this chain. First, we derive the exact population loss of the state-scaled one-flow class, including nonnegative-boundary and degenerate-state cases. Second, an orthogonal net-drift/turnover parameterization gives an exact information decomposition and an explicit representation--estimation criterion; a supplied exactly solvable case reduces the criterion to one scalar crossover. Third, we report a single real-data study in which the two neural classes differ only in flow separation, together with the strongest same-information baseline and a simple operational baseline. The empirical result is deliberately reported as mixed: the structural contrast is favorable on average at 24 hours but does not pass the original confirmatory gate, and neither neural flow class dominates the practical baselines at long horizons.

At fixed $x$, Equation~\eqref{eq:intro_transition} is a varying-coefficient regression \citep{hastie1993}. We do not claim the general construction as novel. The contribution is the exact geometry generated by the basis $(1-y,-y)$ and its implication for choosing between one signed flow and two nonnegative flows. Neural differential models provide flexible parameterizations of unknown dynamics \citep{chen2018}; here the neural networks instantiate the two model classes rather than constitute a separate architectural contribution.
